\chapter{Lua Voxel Manipulators}
\label{ch:advmap/lvm}

Die Funktionen aus dem Kapitel Grundlegende
Kartenoperationen~(Kapitel~\ref{ch:map/environment}) sind bequem und
einfach zu benutzen, für große Bereiche aber ineffizient. Bei jedem
Aufruf von \texttt{set\_\allowbreak{}node} oder
\texttt{get\_\allowbreak{}node} muss dein Mod mit der Engine
kommunizieren. Das führt zu ständigen einzelnen Kopiervorgängen zwischen
Engine und Mod, was langsam ist und die Leistung deines Spiels schnell
verschlechtert. Ein \emph{Lua Voxel Manipulator} (LVM)\index{lvm@LVM}
kann hier die bessere Wahl sein.

\section{Konzept}\label{konzept}

Mit einem LVM lädst du große Bereiche der Karte in den Speicher deines
Mods. Diese Daten kannst du dann lesen und schreiben, ohne weiter mit
der Engine zu kommunizieren und ohne dass
Callbacks\index{callback@Callback} laufen -- diese Operationen sind also
sehr schnell. Danach schreibst du den Bereich zurück in die Engine und
lässt die Beleuchtung berechnen.

\section{Daten in den LVM lesen}\label{daten-in-den-lvm-lesen}

In einen LVM kannst du nur einen Quader laden. Du musst also die
kleinste und größte Position ermitteln, die du verändern willst. Dann
erzeugst du einen LVM und liest den Bereich ein, zum Beispiel so:

\begin{Shaded}
\begin{Highlighting}[]
\KeywordTok{local} \VariableTok{vm} \OperatorTok{=} \VariableTok{core}\OperatorTok{.}\NormalTok{get\_voxel\_manip}\OperatorTok{()}
\KeywordTok{local} \VariableTok{emin}\OperatorTok{,} \VariableTok{emax} \OperatorTok{=} \VariableTok{vm}\OperatorTok{:}\NormalTok{read\_from\_map}\OperatorTok{(}\VariableTok{pos1}\OperatorTok{,} \VariableTok{pos2}\OperatorTok{)}
\end{Highlighting}
\end{Shaded}

Aus Leistungsgründen liest ein LVM fast nie genau den angegebenen
Bereich, sondern meist einen größeren. Diesen größeren Bereich geben
\texttt{emin} und \texttt{emax} an (\emph{emerged min pos} und
\emph{emerged max pos}). Der LVM lädt den Bereich für dich -- egal ob
dafür aus dem Arbeitsspeicher, von der Festplatte gelesen oder der
Kartengenerator\index{kartengenerierung@Kartengenerierung} aufgerufen
werden muss.

\begin{hinweis}{LVMs und Kartengenerierung}

Verwende
\texttt{core.\allowbreak{}get\_\allowbreak{}voxel\_\allowbreak{}manip()}\index{core@\texttt{core}!get_voxel_manip@\texttt{core.\allowbreak{}get\_\allowbreak{}voxel\_\allowbreak{}manip}}
nicht in der Kartengenerierung, das kann zu Fehlern führen. Verwende
dort stattdessen
\texttt{core.\allowbreak{}get\_\allowbreak{}mapgen\_\allowbreak{}object(\textquotedbl{}voxelmanip\textquotedbl{})}\index{core@\texttt{core}!get_mapgen_object@\texttt{core.\allowbreak{}get\_\allowbreak{}mapgen\_\allowbreak{}object}}.

\end{hinweis}

\section{Nodes lesen}\label{nodes-lesen}

Um die Node-Typen an bestimmten Positionen zu lesen, verwendest du
\texttt{get\_\allowbreak{}data()}. Die Methode gibt ein flaches Array
zurück, in dem jeder Eintrag für den Typ eines Nodes steht.

\begin{Shaded}
\begin{Highlighting}[]
\KeywordTok{local} \VariableTok{data} \OperatorTok{=} \VariableTok{vm}\OperatorTok{:}\NormalTok{get\_data}\OperatorTok{()}
\end{Highlighting}
\end{Shaded}

\texttt{param2} und Lichtwerte erhältst du mit
\texttt{get\_\allowbreak{}light\_\allowbreak{}data()} und
\texttt{get\_\allowbreak{}param2\_\allowbreak{}data()}.

Mit \texttt{emin} und \texttt{emax} berechnest du, wo ein Node in diesen
flachen Arrays liegt. Die Hilfsklasse
\texttt{VoxelArea}\index{voxelarea@\texttt{VoxelArea}} übernimmt die
Berechnung für dich.

\begin{Shaded}
\begin{Highlighting}[]
\KeywordTok{local} \VariableTok{a} \OperatorTok{=} \VariableTok{VoxelArea}\OperatorTok{:}\NormalTok{new}\OperatorTok{(\{}
    \VariableTok{MinEdge} \OperatorTok{=} \VariableTok{emin}\OperatorTok{,}
    \VariableTok{MaxEdge} \OperatorTok{=} \VariableTok{emax}\OperatorTok{,}
\OperatorTok{\})}

\CommentTok{{-}{-} Index des Nodes ermitteln}
\KeywordTok{local} \VariableTok{idx} \OperatorTok{=} \VariableTok{a}\OperatorTok{:}\NormalTok{index}\OperatorTok{(}\VariableTok{x}\OperatorTok{,} \VariableTok{y}\OperatorTok{,} \VariableTok{z}\OperatorTok{)}

\CommentTok{{-}{-} Node lesen}
\FunctionTok{print}\OperatorTok{(}\VariableTok{data}\OperatorTok{[}\VariableTok{idx}\OperatorTok{])}
\end{Highlighting}
\end{Shaded}

Dabei fällt auf, dass \texttt{data[idx]} eine ganze Zahl ist. Die Engine
speichert Nodes aus Leistungsgründen nicht als Strings, sondern als
Zahl, die \emph{Content-ID\index{content-id@Content-ID}} heißt. Die
Content-ID eines Node-Typs erhältst du mit
\texttt{get\_\allowbreak{}content\_\allowbreak{}id()}:

\begin{Shaded}
\begin{Highlighting}[]
\KeywordTok{local} \VariableTok{c\_stone} \OperatorTok{=} \VariableTok{core}\OperatorTok{.}\NormalTok{get\_content\_id}\OperatorTok{(}\StringTok{"default:stone"}\OperatorTok{)}
\end{Highlighting}
\end{Shaded}

Dann kannst du prüfen, ob der Node Stein ist:

\begin{Shaded}
\begin{Highlighting}[]
\KeywordTok{local} \VariableTok{idx} \OperatorTok{=} \VariableTok{a}\OperatorTok{:}\NormalTok{index}\OperatorTok{(}\VariableTok{x}\OperatorTok{,} \VariableTok{y}\OperatorTok{,} \VariableTok{z}\OperatorTok{)}
\ControlFlowTok{if} \VariableTok{data}\OperatorTok{[}\VariableTok{idx}\OperatorTok{]} \OperatorTok{==} \VariableTok{c\_stone} \ControlFlowTok{then}
    \FunctionTok{print}\OperatorTok{(}\StringTok{"Ist Stein!"}\OperatorTok{)}
\ControlFlowTok{end}
\end{Highlighting}
\end{Shaded}

Die Content-IDs von Node-Typen können sich während der Ladezeit noch
ändern. Ermittle sie deshalb nicht während dieser Zeit.

Die Nodes im Datenarray eines LVM liegen in umgekehrter
Koordinatenreihenfolge. Iteriere deshalb immer in der Reihenfolge
\texttt{z, y, x}:

\begin{Shaded}
\begin{Highlighting}[]
\ControlFlowTok{for} \VariableTok{z} \OperatorTok{=} \VariableTok{pos1}\OperatorTok{.}\VariableTok{z}\OperatorTok{,} \VariableTok{pos2}\OperatorTok{.}\VariableTok{z} \ControlFlowTok{do}
    \ControlFlowTok{for} \VariableTok{y} \OperatorTok{=} \VariableTok{pos1}\OperatorTok{.}\VariableTok{y}\OperatorTok{,} \VariableTok{pos2}\OperatorTok{.}\VariableTok{y} \ControlFlowTok{do}
        \ControlFlowTok{for} \VariableTok{x} \OperatorTok{=} \VariableTok{pos1}\OperatorTok{.}\VariableTok{x}\OperatorTok{,} \VariableTok{pos2}\OperatorTok{.}\VariableTok{x} \ControlFlowTok{do}
            \CommentTok{{-}{-} vi (voxel index) ist ein üblicher Name}
            \KeywordTok{local} \VariableTok{vi} \OperatorTok{=} \VariableTok{a}\OperatorTok{:}\NormalTok{index}\OperatorTok{(}\VariableTok{x}\OperatorTok{,} \VariableTok{y}\OperatorTok{,} \VariableTok{z}\OperatorTok{)}
            \ControlFlowTok{if} \VariableTok{data}\OperatorTok{[}\VariableTok{vi}\OperatorTok{]} \OperatorTok{==} \VariableTok{c\_stone} \ControlFlowTok{then}
                \FunctionTok{print}\OperatorTok{(}\StringTok{"Ist Stein!"}\OperatorTok{)}
            \ControlFlowTok{end}
        \ControlFlowTok{end}
    \ControlFlowTok{end}
\ControlFlowTok{end}
\end{Highlighting}
\end{Shaded}

Der Grund dafür liegt in der Rechnerarchitektur. Aus dem Arbeitsspeicher
zu lesen, ist recht teuer, deshalb haben Prozessoren mehrere
Cache-Ebenen. Liegen die angeforderten Daten im Cache, sind sie sehr
schnell verfügbar. Liegen sie nicht dort, kommt es zu einem
\emph{Cache-Miss}, und die Daten werden aus dem Arbeitsspeicher geholt.
Dabei werden auch die umliegenden Daten mitgeholt und ersetzen Daten im
Cache, weil das Programm sehr wahrscheinlich bald Daten aus der Nähe
anfordert. Eine gute Optimierungsregel ist daher, Daten der Reihe nach
zu lesen und \emph{Cache-Thrashing} zu vermeiden.

\section{Nodes schreiben}\label{nodes-schreiben}

Zuerst setzt du die neue Content-ID im Datenarray:

\begin{Shaded}
\begin{Highlighting}[]
\KeywordTok{local} \VariableTok{c\_air} \OperatorTok{=} \VariableTok{core}\OperatorTok{.}\NormalTok{get\_content\_id}\OperatorTok{(}\StringTok{"air"}\OperatorTok{)}

\ControlFlowTok{for} \VariableTok{z} \OperatorTok{=} \VariableTok{pos1}\OperatorTok{.}\VariableTok{z}\OperatorTok{,} \VariableTok{pos2}\OperatorTok{.}\VariableTok{z} \ControlFlowTok{do}
    \ControlFlowTok{for} \VariableTok{y} \OperatorTok{=} \VariableTok{pos1}\OperatorTok{.}\VariableTok{y}\OperatorTok{,} \VariableTok{pos2}\OperatorTok{.}\VariableTok{y} \ControlFlowTok{do}
        \ControlFlowTok{for} \VariableTok{x} \OperatorTok{=} \VariableTok{pos1}\OperatorTok{.}\VariableTok{x}\OperatorTok{,} \VariableTok{pos2}\OperatorTok{.}\VariableTok{x} \ControlFlowTok{do}
            \KeywordTok{local} \VariableTok{vi} \OperatorTok{=} \VariableTok{a}\OperatorTok{:}\NormalTok{index}\OperatorTok{(}\VariableTok{x}\OperatorTok{,} \VariableTok{y}\OperatorTok{,} \VariableTok{z}\OperatorTok{)}
            \ControlFlowTok{if} \VariableTok{data}\OperatorTok{[}\VariableTok{vi}\OperatorTok{]} \OperatorTok{==} \VariableTok{c\_stone} \ControlFlowTok{then}
                \VariableTok{data}\OperatorTok{[}\VariableTok{vi}\OperatorTok{]} \OperatorTok{=} \VariableTok{c\_air}
            \ControlFlowTok{end}
        \ControlFlowTok{end}
    \ControlFlowTok{end}
\ControlFlowTok{end}
\end{Highlighting}
\end{Shaded}

Bist du mit dem Setzen fertig, lädst du das Datenarray zurück in die
Engine:

\begin{Shaded}
\begin{Highlighting}[]
\VariableTok{vm}\OperatorTok{:}\NormalTok{set\_data}\OperatorTok{(}\VariableTok{data}\OperatorTok{)}
\VariableTok{vm}\OperatorTok{:}\NormalTok{write\_to\_map}\OperatorTok{(}\KeywordTok{true}\OperatorTok{)}
\end{Highlighting}
\end{Shaded}

Für Licht- und \texttt{param2}-Daten gibt es die passend benannten
Methoden \texttt{set\_\allowbreak{}light\_\allowbreak{}data()} und
\texttt{set\_\allowbreak{}param2\_\allowbreak{}data()}.

\texttt{write\_\allowbreak{}to\_\allowbreak{}map()} erwartet einen
Wahrheitswert: \texttt{true}, wenn die Beleuchtung berechnet werden
soll. Übergibst du \texttt{false}, musst du die Beleuchtung später mit
\texttt{core.\allowbreak{}fix\_\allowbreak{}light}\index{core@\texttt{core}!fix_light@\texttt{core.\allowbreak{}fix\_\allowbreak{}light}}
neu berechnen.

\section{Beispiel}\label{beispiel}

\begin{Shaded}
\begin{Highlighting}[]
\KeywordTok{local} \KeywordTok{function}\NormalTok{ grass\_to\_dirt}\OperatorTok{(}\VariableTok{pos1}\OperatorTok{,} \VariableTok{pos2}\OperatorTok{)}
    \KeywordTok{local} \VariableTok{c\_dirt} \OperatorTok{=} \VariableTok{core}\OperatorTok{.}\NormalTok{get\_content\_id}\OperatorTok{(}\StringTok{"default:dirt"}\OperatorTok{)}
    \KeywordTok{local} \VariableTok{c\_grass} \OperatorTok{=}
        \VariableTok{core}\OperatorTok{.}\NormalTok{get\_content\_id}\OperatorTok{(}\StringTok{"default:dirt\_with\_grass"}\OperatorTok{)}

    \CommentTok{{-}{-} Daten in den LVM lesen}
    \KeywordTok{local} \VariableTok{vm} \OperatorTok{=} \VariableTok{core}\OperatorTok{.}\NormalTok{get\_voxel\_manip}\OperatorTok{()}
    \KeywordTok{local} \VariableTok{emin}\OperatorTok{,} \VariableTok{emax} \OperatorTok{=} \VariableTok{vm}\OperatorTok{:}\NormalTok{read\_from\_map}\OperatorTok{(}\VariableTok{pos1}\OperatorTok{,} \VariableTok{pos2}\OperatorTok{)}
    \KeywordTok{local} \VariableTok{a} \OperatorTok{=} \VariableTok{VoxelArea}\OperatorTok{:}\NormalTok{new}\OperatorTok{(\{}
        \VariableTok{MinEdge} \OperatorTok{=} \VariableTok{emin}\OperatorTok{,}
        \VariableTok{MaxEdge} \OperatorTok{=} \VariableTok{emax}\OperatorTok{,}
    \OperatorTok{\})}
    \KeywordTok{local} \VariableTok{data} \OperatorTok{=} \VariableTok{vm}\OperatorTok{:}\NormalTok{get\_data}\OperatorTok{()}

    \CommentTok{{-}{-} Daten ändern}
    \ControlFlowTok{for} \VariableTok{z} \OperatorTok{=} \VariableTok{pos1}\OperatorTok{.}\VariableTok{z}\OperatorTok{,} \VariableTok{pos2}\OperatorTok{.}\VariableTok{z} \ControlFlowTok{do}
        \ControlFlowTok{for} \VariableTok{y} \OperatorTok{=} \VariableTok{pos1}\OperatorTok{.}\VariableTok{y}\OperatorTok{,} \VariableTok{pos2}\OperatorTok{.}\VariableTok{y} \ControlFlowTok{do}
            \ControlFlowTok{for} \VariableTok{x} \OperatorTok{=} \VariableTok{pos1}\OperatorTok{.}\VariableTok{x}\OperatorTok{,} \VariableTok{pos2}\OperatorTok{.}\VariableTok{x} \ControlFlowTok{do}
                \KeywordTok{local} \VariableTok{vi} \OperatorTok{=} \VariableTok{a}\OperatorTok{:}\NormalTok{index}\OperatorTok{(}\VariableTok{x}\OperatorTok{,} \VariableTok{y}\OperatorTok{,} \VariableTok{z}\OperatorTok{)}
                \ControlFlowTok{if} \VariableTok{data}\OperatorTok{[}\VariableTok{vi}\OperatorTok{]} \OperatorTok{==} \VariableTok{c\_grass} \ControlFlowTok{then}
                    \VariableTok{data}\OperatorTok{[}\VariableTok{vi}\OperatorTok{]} \OperatorTok{=} \VariableTok{c\_dirt}
                \ControlFlowTok{end}
            \ControlFlowTok{end}
        \ControlFlowTok{end}
    \ControlFlowTok{end}

    \CommentTok{{-}{-} Daten schreiben}
    \VariableTok{vm}\OperatorTok{:}\NormalTok{set\_data}\OperatorTok{(}\VariableTok{data}\OperatorTok{)}
    \VariableTok{vm}\OperatorTok{:}\NormalTok{write\_to\_map}\OperatorTok{(}\KeywordTok{true}\OperatorTok{)}
\KeywordTok{end}
\end{Highlighting}
\end{Shaded}

\section{Jetzt bist du dran}\label{jetzt-bist-du-dran}

\begin{itemize}
\tightlist
\item
  Schreibe
  \texttt{replace\_\allowbreak{}in\_\allowbreak{}area(from, to, pos1, pos2)},
  das im angegebenen Bereich jedes Vorkommen von \texttt{from} durch
  \texttt{to} ersetzt; \texttt{from} und \texttt{to} sind Node-Namen.
\item
  Schreibe eine Funktion, die alle Truhen um 90° dreht.
\item
  Schreibe eine Funktion, die mit einem LVM bemoosten Bruchstein auf
  benachbarten Stein und Bruchstein übergreifen lässt. Breitet sich der
  Moosbewuchs bei deiner Umsetzung jedes Mal weiter als einen Node aus?
  Wenn ja, wie könntest du das verhindern?
\end{itemize}
